\documentclass[11pt]{article}
\usepackage[margin=1in]{geometry}
\usepackage{times}
\usepackage{hyperref}
\usepackage{titlesec}
\usepackage{parskip}
\usepackage{enumitem}

\hypersetup{colorlinks=true, urlcolor=blue}
\titleformat{\section}{\normalfont\bfseries}{\thesection.}{0.5em}{}
\setlength{\parskip}{0.6em}
\setlist[enumerate]{nosep, leftmargin=1.5em}

\title{\vspace{-2em}Presentation 1 Reflection: \\ \large Dialing In Espresso}
\author{Sean Balbale \\ CPSC 403: Computer Science Seminar}
\date{October 13, 2026}

\begin{document}
\maketitle
\vspace{-1.5em}

\section{What I intended}
I wanted the audience to leave with one rule: when a shot tastes wrong, check the clock, change only the grind by one step, and taste again. The recipe slide gave them three numbers (18 g in, 36 g out, 25 to 30 seconds), the grind slide tied direction to taste, and the dial-in loop made it something they could run on any machine.

\section{What landed}
My delivery. I was told I was well spoken and talked to the audience instead of reading the slides, which matches how I built the deck: one short line per slide, with the explanation spoken.
\section{What did not land}
Photos were the main criticism. The bottomless portafilter slide used a drawn schematic, but its whole point is that you can see channeling, and a drawing asks the audience to take my word for it. The gear slide names the Gaggia Classic Pro E24, Fellow Opus 2, and bottomless portafilter without showing any of them, even though I own all three.

There was also too much text per slide. People could not finish reading a slide before I moved on, so the slides competed with me instead of backing me up. Slides like the gear slide and the what-happens-in-a-shot slide put a heading, labels, and a sentence under each of several items.

I also ran short, and nothing showed the loop run with actual numbers.

\section{What I will change for Presentation 2}
\begin{enumerate}
  \item Use real photos, video, or screenshots of the actual subject wherever I would otherwise draw a diagram or just name something.
  \item Cut each slide to a heading, one line, and an image, and split any slide that takes more than a few seconds to read.
  \item Fill the runtime with concrete material, such as a worked example with real numbers or a demo, not more text per slide.
  \item Rehearse out loud with a timer against the required runtime and fix the length before the talk.
\end{enumerate}

\end{document}